\documentclass[11pt,reqno]{amsart}
\usepackage[T1]{fontenc}\usepackage[utf8]{inputenc}\usepackage{lmodern}
\usepackage[a4paper,margin=23mm]{geometry}
\usepackage{amsmath,amssymb,mathtools,microtype,array,xurl}
\usepackage[hidelinks]{hyperref}
\newtheorem{theorem}{Theorem}\newtheorem{lemma}[theorem]{Lemma}
\newtheorem{corollary}[theorem]{Corollary}
\title{Finite Reduction of Mixed Square-Source Closed Sets\\for the Erd\H{o}s--Straus Equation}
\author{The Clankers}\date{18 September 2026}
\begin{document}
\begin{abstract}
Square-multiplied shifted factors impose a finite reduction on every closed set
of prescribed cardinality in the external-nonresidue graph, without restricting
its stabilizer index or its vertex types. A set of size $m$ forces $p<16^{m^2}$.
A complete finite cofactor enumeration excludes all sets of size at most five.
An integer carry preserves the full selector even when its square residual overlaps
the base modulus. These statements do not prove universal selector occupancy.
\end{abstract}\maketitle

Let $p$ be prime in $1,121,169,289,361,529\pmod{840}$. Put
\[
 V_p=\{q<p:q\text{ prime},\ (q/p)=-1\},\qquad
 \sigma_q=\begin{cases}1&q\equiv3\pmod4,\\3&q\equiv1\pmod4.\end{cases}
\]
The canonical graph and its factorwise reciprocity rules are inherited [1].
The primes $2,3,5,7$ are residues, so $V_p\ne\varnothing$ and $\min V_p\ge11$.
For odd $t>0$ define the separately marked source
\[
 R=\sigma_q qt^2<3p,\qquad A_{q,t}=(p+R)/4.
\]
It is an integer below $p$, coprime to $pR$, and $(A_{q,t}/p)=-1$.
Every such source therefore has a nonresidue prime factor below $p$, distinct
from $q$. Every factor is an edge, retaining its full valuation.
Call $C\subseteq V_p$ \emph{closed} if all those nonresidue factors remain in $C$
at every $q\in C$ and every valid $t$. This is stronger than canonical closure
and independent of selector occupancy. Previously, only all-$3\bmod4$ triangles
and small primitive components had a proved square-source escape [2].

\begin{lemma}[An original source avoiding a finite vertex set]\label{prefix}
Let $q\in C$, $|C|=m$. Define
\[
 (b_0,\ldots,b_7)=(1,2,3,4,5,6,9,15),\qquad
 b_n=\left\lceil\frac{(2n+9)(3^n-1)}{18}\right\rceil\ (n\ge8).
\]
Among the $b_{m-1}$ integers $F_q(j)=A_{q,2j+1}$, $0\le j<b_{m-1}$,
at least one is coprime to every prime of $C$. If all these sources are below $p$,
that source has an actual nonresidue factor outside $C$.
\end{lemma}
\begin{proof}
Write $F_q(j)=A_{q,1}+\sigma_q qj(j+1)$. At $r\in C\setminus\{q\}$ it has
zero or two roots, since its discriminant is the nonzero residue $-\sigma_q qp$.
The two roots sum to $-1$. The source has no root modulo $q$.
For a subset of the $n\le m-1$ active forbidden primes, CRT gives $2^k$ roots
modulo their product $d$, paired by $a\mapsto d-1-a$ without fixed points.
For one pair, its count in $[0,N)$ differs from $2N/d$ by less than one: after
complete periods are removed, a prefix of length at most $(d-1)/2$ contains at
most one root and a longer proper prefix contains at least one. Inclusion--exclusion
therefore gives the exact count $Z$ with
\[
 \left|Z-N\prod_r(1-2/r)\right|<(3^n-1)/2\quad(n>0).       \tag{1}
\]
The increasing active primes satisfy $r_i\ge2i+9$, so
$\prod_r(1-2/r)\ge9/(2n+9)$. This proves the formula for $n\ge8$.
For $1\le n\le5$, each prime at least 11 covers at most one index in $[0,n]$:
two distinct roots have sum at least 10, while the largest possible sum there is
$2n-1\le9$. For $n=6$, in a prefix of length 9 only 11 and 13 can cover two indices,
so at most 8 are covered. For $n=7$, in a prefix of length 15, 11 and 13 cover at
most 3 each, 17,19,23 at most 2 each, and all others at most 1: the total is at most 14.
For $n=0$ every index works. The constants increase and satisfy $b_n\le4^n$;
for the formula use $(2n+9)/18<(4/3)^n$, proved inductively. Finally use the
negative character of an actual source below $p$ to obtain its outside prime.
\end{proof}

\begin{theorem}[All-cardinality finite reduction]\label{height}
A nonempty closed set $C$ of size $m$ has $m\ge2$. For
$H=(2b_{m-1}-1)^2$, every $q\in C$ satisfies
\[
 \boxed{\sigma_q qH\ge3p,\qquad q\ge p/H,\qquad p<(4H)^m<16^{m^2}.} \tag{2}
\]
It contains a simple canonical cycle of length $2\le l\le m$ with equations
$p+\sigma_iq_i=4c_iq_{i+1}$. Rotate to $c_0=v=\min c_i$. Then
\[
 \boxed{1\le v<(H+3)/4,\quad
 v\le c_i\le\left\lfloor\frac{vH-1}{4v-3}\right\rfloor.}          \tag{3}
\]
Each coefficient word in this finite domain reconstructs at most one candidate
prime cycle, including both original vertex types.
\end{theorem}
\begin{proof}
Lemma~\ref{prefix} proves the first inequality in (2); otherwise all its sources
are valid and one has an outside factor. Since $\sigma_q\le3$, $q\ge p/H$.
Canonical sources always give an outgoing edge, and self-edges are impossible,
so finite closure supplies a simple cycle. Set
$D=4^l\prod c_i$, $S=\prod\sigma_i$. At each cyclic starting point $i$, iterate
\[
 D_0=1,\ T_0=0,\quad
 D_{j+1}=4c_{i+j}D_j,\quad T_{j+1}=\sigma_{i+j}T_j+D_j.
\]
Its final numerator $T_i$ satisfies $(D-S)q_i=pT_i$, and $D>S$.
Therefore $p\mid D-S$, and, since the distinct primes $q_i$ have gcd one,
\[
 h=\gcd_iT_i,\qquad \boxed{p=(D-S)/h,\quad q_i=T_i/h.}            \tag{4}
\]
All integer, prime, character, type, range and edge conditions must be checked
on the reconstructed candidate. Since $c_i=A_{q_i,1}/q_{i+1}<H$,
$p\le D-S<(4H)^m<16^{m^2}$, using $b_{m-1}\le4^{m-1}$.
For the largest cycle vertex $Q$, its incoming edge gives
$4vQ<p+3Q$, so $Q<p/(4v-3)$. Each canonical source is less than $vp/(4v-3)$,
hence $c_i<vH/(4v-3)$. Taking the exact integer endpoint and requiring it at least
$v$ proves (3). The reconstruction (4) is a complete necessary finite reduction,
not a favourable-factor hypothesis.
\end{proof}

\begin{theorem}[No closed set of size at most five]\label{five}
At every hard prime, every nonempty set of at most five external nonresidue primes
has an actual nonresidue factor outside it in some valid source with
$t\in\{1,3,5,7,9\}$.
\end{theorem}
\begin{proof}
At a fixed vertex the exact source gcd is
\[
 \gcd(F_q(i),F_q(j))=\gcd(F_q(i),(j-i)(i+j+1)).                  \tag{5}
\]
For the first five ports its prime factors are at most 8, whereas vertices are at
least 11. Five valid ports therefore need five distinct successors, impossible
inside a set of at most five after excluding the source itself.
Thus closure would give $81\sigma_q q\ge3p$ at every vertex. The finite domain
(3) becomes
\[
 2\le l\le5,\quad\sigma_i\in\{1,3\},\quad1\le v\le20,
 \quad c_0=v,\quad v\le c_i\le\left\lfloor\frac{81v-1}{4v-3}\right\rfloor. \tag{6}
\]
Retain the inequalities $81\sigma_iT_i\ge3(D-S)$ and reconstruct by (4).
The supplied standard-library enumeration and its second implementation exhaust
this domain: 1,381,117,764 rooted profiles are bounded; 5,922,259 pass the range
inequalities; 1,915 reconstruct typed hard-class integers; six have every integer
prime; four rooted profiles have every vertex a nonresidue. Up to rotation these
are exactly
\[
\begin{array}{r|l}
4201&61,137,1153,383,191\\
12601&409,3457,5743,2293,487\\
26209&661,881,7213,5981,5519.
\end{array}                                                   \tag{7}
\]
For completeness, the exact exhaustive reduction used in that computation is as
follows. Fix other positive coefficients and set $c_j=x$. Write
$D-S=D_0x-S$, $T_i=A_ix+B_i$, where $A_i\ge0$, $B_i>0$.
Then
\[
 \frac{d}{dx}\frac{T_i}{D-S}
 =\frac{-A_iS-D_0B_i}{(D_0x-S)^2}<0.                           \tag{8}
\]
With all unassigned coordinates at their lower bound $v$, failure of any range
inequality excludes the entire remaining box. Otherwise the next coordinate
satisfies
$[3D_0-81\sigma_iA_i]x\le81\sigma_iB_i+3S$; each positive bracket gives an
exact floor endpoint. The main checker recurses through every permitted integer.
The second uses binary search of (8) and independently reconstructs the numerators
from $4c_iT_{i+1}=D-S+\sigma_iT_i$. Both produce the same full range-profile hash.
The ledger retains every typed candidate and a proper composite divisor or failed
character when rejected. All trial divisions are exhaustive, not probable-prime tests.

Each cycle in (7) already has five vertices. At its first vertex, multiplier 9 gives
respectively
\[
 1462=2\cdot17\cdot43,\qquad5911=23\cdot257,\qquad11014=2\cdot5507.
\]
The nonresidue primes $43,23,5507$ are outside their respective cycles. Every
factorization, primality, character and size condition is checked. Thus every
remaining candidate leaves the proposed set. This is a finite computer-assisted
proof after the universal reduction, not a scan of primes up to a chosen bound.
\end{proof}

Starting at any external nonresidue prime, a breadth-first exploration of those
five valid ports must reach at least six distinct vertices after at most 25 source
factorizations. Otherwise the explored set of at most five vertices would be
closed. All source integers are below $p$, so ordinary trial division makes this a
terminating factor-supply algorithm. It is not a terminating ES selector.
More generally Theorem~\ref{height} uses only its initial $b_{m-1}$ ports; for
$p\ge16^{m^2}$ they force at least $m+1$ reachable vertices within $m b_{m-1}$
source factorizations.

The argument allows all quotient indices and both vertex types. Opposite mixed
edges cannot be deleted: at $p=12889$, the canonical pair $43\leftrightarrow61$
has sources $53\cdot61$ and $2^2\cdot19\cdot43$. Both selectors are empty at
\emph{all} valid squares there: $t=1,3,\ldots,29$ at 43 and $t=1,3,\ldots,13$ at 61.
The accompanying independent checker exhausts all 246 original divisors across
these 22 sources. Thus outside-factor escape is not immediate local occupancy.
A valid path $43\xrightarrow{3}3319\xrightarrow{1}1013\xrightarrow{1}11$ ends at
$A=3225,u=9$ and returns
\[
 \frac4{12889}=\frac1{3225}+\frac1{1357856150}+\frac1{3789366}.
\]

\paragraph{The full square residual and the integer coefficient.}
For the original source $A=\prod l^{e_l}$ retain each
$\beta_l\in[-e_l,e_l]$ and $u=\prod l^{e_l+\beta_l}$. Set $R_0=\sigma_q q$,
$\eta_E=1$, $\eta_M=p$. For every actual base candidate
$u\mid A^2$, $R_0\mid4u+\eta_c$, keep
\[
 k=\frac{4u+\eta_c}{R_0},\quad d=k\bmod t^2,\quad b=(k-d)/t^2.
\]
Then
\[
 |c_{q,t}|=[Z^0]\sum_{u}Z^{d(u)},\qquad
 u=\frac{R_0(d+t^2b)-\eta_c}{4}.                           \tag{9}
\]
These are integer coefficients with original labels. The fibre of a digit is all
allowed $b$ in (9), not a chosen representative. No coprimality of $R_0,t$ is used.
At $p=1129,q=t=11$, $A=615$ and $u=1845$ passes both mod 11 and mod 121, but fails
mod 1331: its carry is 671, with digit 66 mod 121. Replacing a product by the lcm would
lose an actual valuation constraint. Across two square sources the common divisor
domain is exactly $\operatorname{Div}(\gcd(A,A')^2)$, and the same integer carry is
retained there. A zero augmentation or a zero after coefficient reduction is not
an absent original divisor.

For a full hit, $d_0=\gcd(A,u)$ gives $(h,r,s)=(d_0^2/u,u/d_0,A/d_0)$.
The ordered outputs are
\[
 E:(A,hs\kappa,phr\kappa),\quad\kappa=(pr+s)/R;
 \qquad M:(A,phs\lambda,phr\lambda),\quad\lambda=(r+s)/R.
\]
Their original ordered inverses are $u=A^2/(Ry-pA)$ and $u=pA^2/(Ry-pA)$.
Direct multiplication proves $4/p$; these are the inherited Type I/II coordinates [3].
No universal nonzero target coefficient is inferred from (9) or from source escape.
The remaining issue is simultaneous failure on larger sets, not a missing
formal inverse of a reduced coefficient.

\paragraph{Reproducibility and scope.}
Run the following commands with the accompanying Python files:
{\small\begin{verbatim}
python cofactor_check.py \
  --out certificates/cofactor_certificate.json
python check_independent.py \
  --input certificates --out independent.json
\end{verbatim}}
The portable example checker separately verifies the 22 sources and their positive
return. Optimized Python is supported. The full workbench adds the complete root
counts, prime-power lifts, mixed reciprocity table and bounded graph replay.
No ES resolution, least-counterexample assertion, historical-priority determination,
Lean build or independent mathematical review is claimed.

\begin{thebibliography}{3}
\bibitem{Bryan} Bryan, C. (2026). \emph{External-nonresidue factor cycle}, \S\S1--7.
CENTL, commit \texttt{42508684c4386b575d7b52e763565345a76bcb9e},
\url{https://github.com/chasebryan/centl/blob/42508684c4386b575d7b52e763565345a76bcb9e/research/erdos-straus/EXTERNAL-NR-FACTOR-CYCLE.md}.
\bibitem{Turn3} The Clankers. (2026). \emph{Primitive sextic cycles and square-source
escape for the Erd\H{o}s--Straus selectors}, \S\S6--7. Commit
\texttt{852f86e085c7bfbb381b85cafb222a5802194e63},
\url{https://github.com/KokunoYumeto/erdos-straus-foundation/pull/18}.
\bibitem{ET} Elsholtz, C., \& Tao, T. (2013). Counting the number of solutions to the
Erd\H{o}s--Straus equation on unit fractions. \emph{Journal of the Australian
Mathematical Society, 94}(1), 50--105. \S2, Propositions2.2 and2.6.
\url{https://arxiv.org/abs/1107.1010}.
\end{thebibliography}
\end{document}
